\section{Finite-volume voxel-spacing-aware texture matrices}

\subsection{Rationale}

The final implementation uses a finite-volume, zero-order-hold representation for selected texture families when
\texttt{weightingNorm = voxelSpacing} is explicitly enabled. The purpose is to preserve the native acquired signal
while representing anisotropic voxels on an isotropic physical lattice before matrix construction. This separates
geometric correction from interpolation-induced signal modification: voxel intensities are not smoothed, mixed, or
interpolated, but an anisotropic voxel is represented by repeated subcells with the same discretized gray level.

Let the image-array spacing be
\[
s = (\Delta_z,\Delta_y,\Delta_x),
\]
and define the finest spacing unit
\[
h = \min(\Delta_z,\Delta_y,\Delta_x).
\]
For each axis, the finite-volume expansion factor is
\[
r_k = \mathrm{round}\left(\frac{\Delta_k}{h}\right), \qquad k \in \{z,y,x\}.
\]
The expanded image $\tilde I$ and mask $\tilde M$ are obtained by zero-order repetition of each original voxel along
axis $k$ by $r_k$ subcells. For the synthetic anisotropic experiment with SimpleITK spacing
\[
(\Delta_x,\Delta_y,\Delta_z)=(1,1,5)\,\mathrm{mm},
\]
the array-order spacing is $(5,1,1)$ and the expansion factor is $(5,1,1)$.

\subsection{Feature-family behavior}

\paragraph{GLRLM.}
The gray-level run length matrix is computed using the standard PyRadiomics backend on the finite-volume
representation $\tilde I$. Run length remains an integer count, but it now counts subcells on the isotropic physical
lattice rather than native anisotropic index voxels. Therefore, a continuous run through a 5-mm slice direction is
represented by five 1-mm subcell steps. Diagonal directions are handled by the standard GLRLM angle generation after
finite-volume expansion; no additional radius rule or angular weighting is introduced.

\paragraph{GLDM.}
The gray-level dependence matrix is likewise computed by the standard GLDM backend on $\tilde I$. Dependence size
remains an integer count of neighboring subcells satisfying the gray-level dependence criterion. This avoids the
previous weighted-dependence construction, which introduced a continuous dependence quantity that then had to be
mapped back to a discrete matrix axis.

\paragraph{GLSZM.}
The gray-level size zone matrix is computed by the standard connected-component algorithm on $\tilde I$. Zone size
remains an integer count, but the count is now measured in isotropic subcells. This is physically preferable to
leaving GLSZM entirely insensitive to voxel anisotropy, because a connected zone spanning thick slices represents a
larger physical volume than the same native voxel count in an isotropic image. Importantly, topology and connectivity
rules remain the standard GLSZM rules; the modification affects the finite-volume representation used before zone
matrix construction, not the feature formulas.

\subsection{Implementation summary}

The implementation was applied to \texttt{radiomics\_modificado\_corregido}. The following behavior is enforced:
\begin{itemize}
  \item If \texttt{weightingNorm} is disabled, the modified library follows standard PyRadiomics behavior.
  \item If \texttt{weightingNorm = voxelSpacing} and spacing is isotropic, the extraction reduces exactly to standard
  behavior.
  \item If \texttt{weightingNorm = voxelSpacing} and spacing is anisotropic, GLRLM, GLDM and GLSZM matrices are
  computed after finite-volume zero-order-hold expansion.
  \item Voxel-based extraction for these finite-volume modes is intentionally not enabled in the current prototype.
\end{itemize}

\subsection{Synthetic validation}

The validation used a deterministic synthetic 3D image with a masked cuboid/spherical region, a smooth in-plane
gradient, through-plane alternating bands and fixed random perturbation. Texture features were extracted with
resampling disabled and original image type only. The compared texture families were GLCM, GLRLM, GLDM, GLSZM and
NGTDM. The experiment used isolated Python environments for the original and modified libraries, with separately
compiled C extensions.

\begin{table}[htbp]
\centering
\caption{Environment and import-path validation.}
\begin{tabular}{lll}
\hline
Library & Python / platform & Imported radiomics path \\
\hline
Original & Python 3.10.9, macOS arm64 &
\texttt{build\_sources/original/radiomics/\_\_init\_\_.py} \\
Modified corrected & Python 3.10.9, macOS arm64 &
\texttt{build\_sources/modified/radiomics/\_\_init\_\_.py} \\
\hline
\end{tabular}
\end{table}

Both libraries imported their own compiled extensions:
\texttt{\_cmatrices.cpython-310-darwin.so} and \texttt{\_cshape.cpython-310-darwin.so}. This confirms that the
comparison did not mix original and modified PyRadiomics imports.

\begin{table}[htbp]
\centering
\caption{Backward compatibility check: original PyRadiomics versus modified library in default mode.}
\begin{tabular}{lrrrr}
\hline
Comparison & Features & Mean absolute difference & Maximum absolute difference & Changed features \\
\hline
Original anisotropic vs modified default anisotropic & 75 & 0.000000 & 0.000000 & 0 \\
\hline
\end{tabular}
\end{table}

\begin{table}[htbp]
\centering
\caption{Activation of \texttt{weightingNorm = voxelSpacing} under anisotropic spacing.}
\begin{tabular}{lrrrrr}
\hline
Family & Features & Changed & Changed at $10^{-5}$ & Median absolute difference & Median relative difference \\
\hline
GLCM  & 24 & 24 & 24 & 0.095800 & 0.101046 \\
GLDM  & 14 & 11 & 11 & 0.117819 & 0.491003 \\
GLRLM & 16 & 16 & 16 & 0.468669 & 0.260970 \\
GLSZM & 16 &  8 &  8 & 0.002072 & 0.400000 \\
NGTDM &  5 &  5 &  5 & 0.022053 & 0.335726 \\
\hline
\end{tabular}
\end{table}

\begin{table}[htbp]
\centering
\caption{Isotropic sanity check: modified default versus modified \texttt{voxelSpacing} on isotropic spacing.}
\begin{tabular}{lrrrr}
\hline
Comparison & Features & Mean absolute difference & Maximum absolute difference & Changed features \\
\hline
Modified default isotropic vs modified voxelSpacing isotropic & 75 & 0.000000 & 0.000000 & 0 \\
\hline
\end{tabular}
\end{table}

\subsection{Interpretation}

The finite-volume implementation preserves backward compatibility exactly in default mode and preserves exact
isotropic equivalence when \texttt{weightingNorm = voxelSpacing}. Under anisotropic spacing, the expected families
activate. GLRLM changes in all evaluated features because run lengths are now measured on the finite-volume lattice.
GLDM changes primarily in dependence-size-sensitive features, while purely gray-level terms remain unchanged. GLSZM
changes selectively: features based on zone size or zone percentage change, whereas features depending only on the
number or distribution of gray levels and zones remain unchanged. This selectivity is physically coherent and supports
the interpretation that the implementation modifies geometric representation rather than globally perturbing all
biomarkers.

\subsection{Computational implications}

The principal computational cost is the finite-volume expansion factor
\[
R = r_z r_y r_x.
\]
For spacing $(5,1,1)$ in array order, $R=5$. Thus, memory use for the temporary image and mask and the computational
load of matrix construction increase approximately five-fold before considering feature-family-specific matrix axes.
GLRLM may also require a larger maximum run-length axis, and GLSZM may require a larger maximum zone-size axis because
the expanded mask contains more subcells. For typical anisotropic CT or MRI with one thick through-plane axis, this is
usually a linear increase of roughly $\Delta_\mathrm{thick}/\Delta_\mathrm{thin}$. Highly anisotropic or multi-axis
non-isotropic data can become substantially heavier, and production code should include a clear warning or guard when
$R$ is large.

\subsection{Manuscript-ready wording}

In the final voxel-spacing-aware mode, GLRLM, GLDM and GLSZM were computed after representing the native anisotropic
voxel grid as an isotropic finite-volume lattice using zero-order hold. This operation preserves native gray-level
values and avoids interpolation-induced smoothing. The subsequent texture matrices are then computed using the
standard PyRadiomics matrix definitions on the expanded lattice. Consequently, run length, dependence size and zone
size remain integer-valued matrix axes, but they are expressed in isotropic physical subcells rather than anisotropic
index voxels. This design preserves standard descriptor semantics more closely than fractional distance weighting,
while allowing anisotropic voxel geometry to influence neighborhood-, run- and zone-based texture descriptors.

\subsection{Limitations}

This synthetic experiment verifies software behavior and activation of spacing-aware execution paths under controlled
conditions. It does not assess clinical performance. The finite-volume approach is conservative with respect to matrix
formulas, but it is still not the unmodified IBSI/PyRadiomics reference computation because matrix construction is
performed on an expanded representation. Non-integer spacing ratios are rounded to integer repeat factors in the
current prototype. Future pull-request preparation should document this explicitly, add unit tests for isotropic
equivalence and anisotropic activation, and define a practical maximum expansion factor for large volumes.

